\section{Introduction}

The representation of numbers in large language models (LLMs) fundamentally tests the boundary between linguistic surface forms and abstract mathematical reasoning. Given the ability of LLMs to solve complex reasoning problems, understanding how they internally map linguistic structures to numeric values is critical. The French counting system presents a classic example of linguistic opacity. While most modern Indo-European languages employ a straightforward base-10 (decimal) system, Standard French retains vigesimal (base-20) elements for numbers between 70 and 99. For instance, the number 80 is expressed as \textit{quatre-vingts}, literally translating to ``four twenties,'' and 97 is \textit{quatre-vingt-dix-sept}, or ``four-twenty-ten-seven.'' This opaque structure raises a fundamental question: does a model's understanding of a number tie directly to its linguistic complexity, or does the model possess an abstract numeric representation robust to such surface-level variations?

There is growing interest in how compositional numeral systems are processed by models. Existing work has shown that earlier language models struggle with French numeral magnitude comparison due to this vigesimal opacity \citep{johnson2020probing}. However, it also becomes increasingly challenging to disentangle the inherent complexity of the French language from the specific mathematical structure of its numbers. On one hand, studies treat ``French'' as a monolithic entity, often comparing it to high-resource decimal languages like English or completely regular systems like Japanese. On the other hand, this ignores natural controlled experiments provided by regional dialectal variations. 

To isolate the cost of this structural complexity---the ``vigesimal tax''---we propose a systematic comparative analysis within the French language itself. We focus on the contrast between \standardfrench (vigesimal) and \swissfrench (decimal) numeral systems. Swiss French utilizes transparent decimal forms for the 70-99 range, such as \textit{septante} (70), \textit{huitante} (80), and \textit{nonante} (90). By holding the underlying language constant and varying only the numeral structure, we can directly observe how the vigesimal structure impacts tokenization fragmentation, decoding accuracy, and magnitude reasoning in state-of-the-art models. 

Our investigation across \gpt, \claude, and \llama reveals that while the ``vigesimal tax'' no longer limits simple decoding, it profoundly disrupts deeper mathematical reasoning. Specifically, we find a 2.77 token increase in fragmentation for \standardfrench numerals compared to \swissfrench numerals. Surprisingly, modern models achieve nearly 100\% accuracy in translating both systems to digits. Yet, when tasked with comparing magnitudes across the two systems, accuracy drops sharply to 76\% for \claude and \llama. This exposes a systemic misalignment in how these models represent equivalent abstract values across dialects.

In summary, our key contributions are:
\begin{itemize}[leftmargin=*,itemsep=0pt,topsep=0pt]
    \item We propose a cross-dialectal framework comparing \standardfrench and \swissfrench to isolate the structural complexity of vigesimal counting.
    \item We quantify the tokenization fragmentation of vigesimal forms, demonstrating a massive token count spike (up to 2.8$\times$) in the 70-99 range across modern tokenizers.
    \item We evaluate zero-shot decoding and magnitude comparison, showing that while modern LLMs decode perfectly, they exhibit significant reasoning gaps during cross-system comparisons.
    \item We uncover a novel ``Standard Overestimation'' bias in \claude, which incorrectly conflates longer textual vigesimal representations with larger numeric magnitudes.
\end{itemize}